\slajdid{02-s011}
\begin{frame}[label=02-s011]
\gusttekst
\begin{columns}[c,onlytextwidth]
\begin{column}{.74\textwidth}
Neka je na primer zahtevima za digitalni sistem definisano
\begin{enumerate}
\item Digitalni sistem ima tri ulazna digitalna signala i jedan izlazni digitalni signal
\item Izlazni digitalni signal je na nivou logičke jedinice samo ako su dva ulaza digitalnih signala na nivou logičke jedinice.
\item U svim ostalim slučajevima izlazni digitalni signal je na nivou logičke nule
\end{enumerate}
Prvi korak je da ulaznim i izlaznim digitalnim signalima dodelimo logičke promenljive. Ulaznim digitalnim signalima ćemo dodeliti logičke promenljive C, B i A, dok ćemo izlaznom digitalnom signalu dodeliti logičku promenljivu F. Radićemo u pozitivnoj logici tako da kada je na primer nivo ulaznog signal jednak logičkoj jedinici odgovarajuća logička promenjiva ima vrednost logičke jedinice.

Na osnovu zahteva koji su postavljeni pred digitalni sistem možemo formirati funkcionalnu tabelu.
\end{column}
\begin{column}{.23\textwidth}
\input{tabele/s011_funkcionalna_tabela.tex}
\end{column}
\end{columns}
\end{frame}
